\documentclass[11pt,a4paper]{article}
\usepackage[utf8]{inputenc}
\usepackage[T1]{fontenc}
\usepackage{amsmath,amssymb}
\usepackage{graphicx}
\usepackage{hyperref}
\usepackage[numbers]{natbib}
\usepackage[margin=1in]{geometry}

\title{Daily Paper Review: EvoMemState ---\\Evolving Memory State for Personal Agents}
\author{Auto-generated Report}
\date{November 6, 2026}

\begin{document}
\maketitle

\section{Paper Summary}

\subsection{Problem and Motivation}
Current personal agent memory systems (MemGPT, Mem0, A-MEM, LightMem, Supermemory) operate under a ``fact-centric paradigm''---they extract, store, update, and retrieve factual records independently. Xiang et al.~\cite{xiang2026evomemstate} identify two fundamental failures of this approach through a preliminary study on MemSyco-Bench. First, fact-centric systems fail to capture inter-memory relations: MemoryBank, LightMem, and Supermemory fail to retrieve all relevant information in 58--62\% of cases, and when all required memories \emph{are} retrieved, answer accuracy is 5.9--33.3\% higher than when some are missing. Second, complete retrieval does not guarantee correct memory use: a fixed self-reflection prompting strategy \emph{reduces} average accuracy from 50.5\% to 48.9\%, offering no consistent improvement. The authors argue for a paradigm shift from fact-centric to \textbf{state-centric} memory, where stored facts are augmented with evolving states that make validity, scope, and relations explicit, and where memory application strategies are learned from experience via reinforcement learning.

\subsection{Method}
EvoMemState introduces two complementary state types that augment the underlying fact store.

\paragraph{Relational States.}
For each persistent fact $f_i$ in the memory store $\mathcal{F}_t$, a relational state captures its evolving context:
\begin{equation}
  r_i^t = (f_i,\; a_i^t,\; \rho_i^t), \quad \mathcal{R}_t = \{r_i^t \mid f_i \in \mathcal{F}_t\},
\end{equation}
where $a_i^t$ specifies the fact's current validity and scope, and $\rho_i^t$ records typed semantic relations to other facts (\textsc{Supersedes}, \textsc{Conflicts\_With}). These states are persistent and incrementally updated as new information arrives, without overwriting original records.

The relational states are organized via an \emph{ontology graph} $\mathcal{G}^{(i)} = (\mathcal{V}_C^{(i)} \cup \mathcal{V}_F^{(i)},\; \mathcal{E}_H^{(i)} \cup \mathcal{E}_R^{(i)})$, where $\mathcal{V}_C$ are concept nodes, $\mathcal{V}_F$ are fact nodes, $\mathcal{E}_H$ are structural edges (\textsc{Parent\_Of}, \textsc{Contains}), and $\mathcal{E}_R$ are relational edges. The graph is constructed iteratively: for each new fact $f_i$, the system decides among \textsc{Create} (new concept node), \textsc{Merge} (attach to existing concept with potential \textsc{Supersedes}/\textsc{Conflicts\_With} edges), or \textsc{Split} (reorganize concept hierarchy). Ontology construction is performed offline using DeepSeek-V4-Flash.

\paragraph{Reasoning States.}
For each query $q_t$, an ephemeral reasoning state is generated by a learned policy:
\begin{equation}
  Z_{q_t} = \Psi(q_t, \mathcal{F}_t, \mathcal{R}_t),
\end{equation}
specifying the applicability, role, and priority of memory evidence for the current request. This is generated as natural language (max 600 characters, 8 lines) by a Qwen3-4B-Instruct model trained via reinforcement learning.

\paragraph{RL-based policy training.}
The Memory State Evolution Policy is trained via Group Relative Policy Optimization (GRPO). For each training query, $G{=}8$ candidate reasoning states are sampled from the current policy. The reward measures marginal improvement:
\begin{equation}
  R_i = \mathrm{Score}\bigl(\mathcal{A}(q, M_q, Z_q^{(i)}),\; y^*\bigr) - \mathrm{Score}\bigl(\mathcal{A}(q, M_q),\; y^*\bigr),
\end{equation}
where $\mathcal{A}$ is a \textbf{frozen} response agent (Qwen3-8B or GPT-4o-mini), $M_q$ is retrieved memory evidence, and $y^*$ is the gold answer. The scoring function is task-specific: LLM-as-judge (G-Eval) for quality scores and token-level F1 for factual accuracy. The GRPO loss is:
\begin{equation}
  \mathcal{L}_{\mathrm{GRPO}}(\theta) = -\frac{1}{G}\sum_{i=1}^{G}\frac{1}{T_i}\sum_{j=1}^{T_i}\min\!\bigl(\rho_{i,j}\,\hat{A}_i,\;\mathrm{clip}(\rho_{i,j}, 1{-}\epsilon, 1{+}\epsilon)\,\hat{A}_i\bigr) + \beta_{\mathrm{KL}}\,D_{\mathrm{KL}}(\pi_\theta \| \pi_{\mathrm{ref}}),
\end{equation}
with clipping $\epsilon{=}0.2$, $\beta_{\mathrm{KL}}{=}0.001$, learning rate $10^{-6}$, rollout batch size~512, trained for 1~epoch on 2~GPUs.

\paragraph{State-aware inference.}
At inference, $K$ semantically similar facts are retrieved as seeds, then \emph{bidirectional graph traversal} follows \textsc{Supersedes} and \textsc{Conflicts\_With} edges to recover connected updates and conflicts. The complete package---retrieved facts, relational states, and the policy-generated reasoning state---is passed to the frozen response agent.

\subsection{Results}
EvoMemState (with Qwen3-8B) achieves \textbf{59.9\% overall average} across four benchmarks, consistently outperforming all 9~baseline systems:

\begin{itemize}
  \item \textbf{MemSyco-Bench} (preference updates, sycophancy): 57.6\% average accuracy vs.\ 45.4\% for the best baseline (A-MEM), a \textbf{+12.2pp} improvement. Sycophancy reduced from 49.1\% to 38.0\%.
  \item \textbf{MemTrapBench} (cognitive traps): 79.1\% average vs.\ 71.0\% for PGMem (\textbf{+8.1pp}). Strongest on Boundaries (91.4\%) and Cognitive Bias (88.8\%).
  \item \textbf{LoCoMo-Refined} (long-term conversational memory): 57.4\% judge score vs.\ 51.3\% for SimpleMem (\textbf{+6.1pp}), 45.5\% token F1.
  \item \textbf{Retrieval completeness}: 85.6\% complete retrieval rate via relation traversal. Conditioned on complete evidence, accuracy improves by +11.5pp on LoCoMo-Refined and +31.2pp on MemSyco-Bench vs.\ baselines---demonstrating that EvoMemState dramatically improves memory \emph{utilization}, not just retrieval.
\end{itemize}

With GPT-4o-mini as the response agent, EvoMemState achieves 61.2\% overall average, confirming the policy transfers across response models.

\paragraph{Cross-system transfer.}
The trained policy can be plugged into other memory systems as a drop-in enhancement, improving NaiveRAG by +4.9\%, MemoryBank by +4.8\%, and A-MEM by +3.4\%. This confirms the policy learns general memory-use reasoning, not system-specific strategies.

\paragraph{Ablations.}
Removing Relational States degrades the overall average by $-4.5$pp ($-8.9$pp on MemSyco-Bench, which specifically tests preference updates). Removing Reasoning States degrades by $-6.9$pp ($-15.0$pp on MemTrapBench, which tests cognitive traps). Removing RL training (using the base Qwen3-4B without GRPO) degrades by $-6.7$pp ($-16.4$pp on MemTrapBench), confirming that the memory-use strategy must be \emph{learned} from downstream feedback, not prompted. Both state types are complementary, degrading on different benchmarks when removed.

\subsection{Limitations}
The ontology is constructed offline using DeepSeek-V4-Flash, preventing real-time memory updates during conversation sessions. The quality of \textsc{Supersedes}/\textsc{Conflicts\_With} relation detection depends entirely on the construction LLM with no error analysis or formal guarantees provided. Only 1~epoch of GRPO training is performed; whether longer training helps is unknown. The 4B policy model is relatively small with no scaling experiments. All benchmarks focus on conversational dialogue memory---no evaluation on task-oriented agents (tool use, web browsing, coding), multi-user scenarios, or truly long-horizon deployments spanning weeks or months. The system is text-only with no multimodal support. No latency or API cost analysis is reported despite the multi-component pipeline. No user study with real users over real interaction histories is conducted.

\subsection{Relevance to My Research}
This paper makes a significant contribution to \textbf{continual learning for LLM agents} (Topic~3). The core insight---that \emph{how} memory is used matters as much as \emph{what} memory is stored---reframes agent memory from a retrieval problem to a policy learning problem. The two-state design elegantly separates persistent knowledge management (Relational States via ontology graphs with typed semantic edges) from adaptive reasoning strategy (Reasoning States via RL-trained policy). This directly addresses catastrophic forgetting through a non-parametric approach: model weights remain frozen while external memory preserves full historical traceability via \textsc{Supersedes} edges rather than overwriting. The RL-based policy training with marginal improvement rewards provides a principled signal for learning memory-use strategies, and the cross-system transfer result (+0.2 to +4.9\% across 6~systems) demonstrates that these strategies are generalizable---a key requirement for lifelong learning agents. The +31.2pp accuracy improvement conditioned on complete evidence (MemSyco-Bench) quantifies the severity of the memory utilization gap in current systems.

\section{Follow-up Research Directions}

\subsection{Direction 1: Online Ontology Evolution with Streaming Memory Updates}
\textbf{Gap:} EvoMemState's ontology construction is performed offline using DeepSeek-V4-Flash, preventing real-time memory updates during conversation sessions. For truly interactive personal agents, memory should evolve incrementally as each new interaction provides information, not in batch after all conversations complete.

\textbf{Approach:} Design a streaming ontology update mechanism that processes new facts immediately upon extraction. Maintain a lightweight \emph{fact buffer} during conversation; after each user turn, extract candidate facts via the response model itself (no external API call) and apply incremental \textsc{Create}/\textsc{Merge}/\textsc{Split} operations using a distilled 1B ontology update model trained to mimic DeepSeek-V4-Flash's decisions. The key challenge is maintaining ontology consistency during rapid updates: implement an \emph{optimistic concurrency} protocol where updates are applied tentatively and validated asynchronously. When a conflict is detected (e.g., a \textsc{Merge} that should have been a \textsc{Split}), a correction is queued and applied during the next idle period. Train the distilled model on (fact, ontology-state, decision) triples collected from the offline pipeline. Evaluate via a ``conversation replay'' protocol: feed multi-session dialogues turn-by-turn and measure ontology quality (precision/recall of relations) and downstream task accuracy compared to the offline oracle.

\textbf{Feasibility:} High. Knowledge distillation from large to small models for structured decisions is well-established. The 1B model adds minimal latency ($<$50ms per update). The main risk is error accumulation over hundreds of incremental updates---the asynchronous validation mechanism should bound this, but empirical testing with long interaction histories (1000+ turns) is needed.

\subsection{Direction 2: Multi-Granularity Reasoning States via Hierarchical RL}
\textbf{Gap:} EvoMemState generates a single flat reasoning state (max 600 characters) per query, regardless of the complexity of the memory-use decision. Simple queries (``What's my preferred coffee?'') need minimal guidance, while complex queries involving multiple conflicting memories, temporal dependencies, and implicit constraints require detailed multi-step reasoning about memory applicability. A fixed-length, single-granularity reasoning state cannot optimally serve both cases.

\textbf{Approach:} Replace the single policy with a \emph{hierarchical} two-level policy. The high-level policy (trained via GRPO) decides the reasoning \emph{strategy type}: (a)~\textsc{Direct} (no reasoning state needed, pass memory directly), (b)~\textsc{Brief} (one-sentence guidance), (c)~\textsc{Detailed} (full multi-step reasoning with conflict resolution, priority assignment, and applicability analysis). The low-level policy generates the actual reasoning state conditioned on the selected strategy. The hierarchical reward decomposes into: high-level reward = accuracy improvement per token of reasoning state (encouraging parsimony), low-level reward = standard marginal improvement. This creates an efficiency-quality tradeoff: the system learns to invest reasoning effort only when the memory situation warrants it. Train via hierarchical GRPO where the high-level and low-level policies alternate updates. Evaluate by measuring accuracy-vs-reasoning-length Pareto curves against the flat baseline.

\textbf{Feasibility:} Moderate. Hierarchical RL adds training complexity (two coupled policies) and requires careful reward shaping to avoid mode collapse (always selecting \textsc{Detailed} or always \textsc{Direct}). The per-token efficiency reward should prevent this, but tuning the tradeoff coefficient requires validation data. The 4B model may have insufficient capacity for nuanced strategy selection; a 7B policy model may be needed.

\subsection{Direction 3: Causal Memory Attribution for Robust Personalization}
\textbf{Gap:} EvoMemState's RL reward measures the marginal improvement of the reasoning state over using memory evidence alone. However, this reward conflates the quality of the reasoning state with confounding factors: the difficulty of the query, the completeness of retrieved evidence, and the frozen agent's baseline capability. A reasoning state that happens to ``unlock'' an easy query gets the same reward as one that resolves a genuinely difficult memory conflict. This can lead to policies that are overfit to easy cases and underperform on the hard ones that matter most.

\textbf{Approach:} Introduce a \emph{causal memory attribution} framework inspired by R$^2$Orch's randomized context design. For each training query, generate reasoning states under two conditions: (a)~with the full ontology (relational states + retrieved facts), and (b)~with a \emph{scrambled} ontology (relations randomly permuted, preserving fact content). The causal effect of the reasoning state is estimated as the difference in accuracy between conditions, isolating the policy's ability to \emph{correctly leverage relational structure} from its ability to simply rephrase facts. Define the causal reward:
\begin{equation}
  R_i^{\mathrm{causal}} = \bigl[\mathrm{Score}(\text{full ontology}, Z_q^{(i)}) - \mathrm{Score}(\text{full ontology})\bigr] - \bigl[\mathrm{Score}(\text{scrambled}, Z_q^{(i)}) - \mathrm{Score}(\text{scrambled})\bigr].
\end{equation}
This reward is positive only when the reasoning state provides value \emph{specifically because} it correctly interprets relational structure, not merely because it reformulates facts. Train with this causal reward using the same GRPO framework. Evaluate whether the causal policy shows improved robustness on out-of-distribution memory configurations (new relation types, deeper ontology graphs, more complex conflict patterns).

\textbf{Feasibility:} Moderate. The scrambled-ontology condition doubles the rollout cost per training query (16 rollouts instead of 8). The causal reward has higher variance than the standard marginal reward, potentially requiring larger rollout groups ($G{=}16$ or $G{=}32$) for stable training. The main risk is that the scrambled condition may be too easy to distinguish from the real one, collapsing the causal signal. Careful design of the scrambling procedure (preserving local structure while disrupting global relations) is essential.

\section{Conclusion}
Xiang et al.\ introduce EvoMemState, a state-centric memory framework that augments personal agent memory with two complementary state types: persistent Relational States capturing inter-memory dependencies via an ontology graph with typed semantic edges, and ephemeral Reasoning States generated by an RL-trained policy specifying how to apply memory evidence to each query. The system achieves 59.9\% overall average across four benchmarks, outperforming all 9~baselines by up to +12.2pp (MemSyco-Bench). The ablations demonstrate that both state types are essential and complementary, with RL training being critical ($-16.4$pp on MemTrapBench without it). The cross-system transfer result---improving 6~different memory systems by up to +4.9\%---confirms that memory-use strategies are generalizable. The most striking finding is the +31.2pp accuracy improvement conditioned on complete evidence, quantifying how severely current systems underutilize available memory. EvoMemState demonstrates that the bottleneck in personal agent memory is not storage or retrieval but \emph{reasoning about when and how to apply stored knowledge}---a problem that can be effectively addressed through learned policies rather than fixed heuristics.

\bibliographystyle{plainnat}
\begin{thebibliography}{1}
\bibitem[Xiang et al.(2026)]{xiang2026evomemstate}
Zhishang Xiang, Yunbo Tang, Zhongqiao Liu, Junkai Chen, Zerui Chen, Qingguo Hu, Yuxuan Zhang, William Grant, Qinggang Zhang, and Jinsong Su.
\newblock {EvoMemState}: Evolving Memory State for Personal Agents.
\newblock \emph{arXiv preprint arXiv:2611.03743}, 2026.
\end{thebibliography}

\end{document}
